% Part: second-order-logic
% Chapter: syntax-and-semantics
% Section: size-of-domain

\documentclass[../../../include/open-logic-section]{subfiles}

\begin{document}

\olfileid{sol}{syn}{siz}

\olsection{Njlentrehake Domain Tanpa Wates lan \usetoken{S}{enumerable}
  \usetoken{P}{domain}}

Himpunan~$M$ tanpa wates (miturut Dedekind) yen lan mung yen
ana fungsi !!a{injective}~$f\colon M \to M$ kang ora
!!{surjective}, yaiku $\ran{f} \neq M$. Ing logika orde
kapisan, kita bisa nimbang !!{function} siji-papan~$f$
lan nyatakake yen fungsi~$\Assign{f}{M}$ kang diwenehake
marang simbol iku ing !!a{structure}~$\Struct{M}$ iku
!!{injective} lan $\ran{f} \neq \Domain{M}$:
\[
\lforall[x][\lforall[y][(\eq[f(x)][f(y)] \to \eq[x][y])]] \land
\lexists[y][\lforall[x][\eq/[y][f(x)]]].
\]
Yen $\Struct{M}$ nyukupi !!{sentence} iki,
$\Assign{f}{M}: \Domain{M} \to \Domain{M}$ iku
!!{injective}, mula $\Domain{M}$ mesthi tanpa wates. Yen
$\Domain{M}$ tanpa wates, mula fungsi kaya mangkono ana,
lan kita bisa netepake $\Assign{f}{M}$ minangka fungsi
iku supaya $\Struct{M}$ nyukupi !!{sentence} kasebut.
Nanging iki mbutuhake basa kita ngemot simbol nonlogis~$f$
kanggo tujuan iku. Ing logika orde kapindho, kita cukup
nyatakake yen fungsi kaya mangkono \emph{ana}. Kita ora
perlu~$f$ maneh, lan oleh !!{sentence} logika orde
kapindho murni
\[
\fn{Inf} \ident \lexists[u][(\lforall[x][\lforall[y][(\eq[u(x)][u(y)]
      \lif \eq[x][y])]] \land \lexists[y][\lforall[x][\eq/[y][u(x)]]])].
\]
$\Sat{M}{\fn{Inf}}$ yen lan mung yen $\Domain{M}$ tanpa
wates. Sabanjure kita bisa netepake
$\fn{Fin} \ident \lnot \fn{Inf}$;
$\Sat{M}{\fn{Fin}}$ yen lan mung yen $\Domain{M}$ winates.
Ora ana siji !!{sentence} logika orde kapisan murni kang bisa
nyatakake yen !!{domain} iku tanpa wates, sanadyan himpunan
tanpa wates saka ukara-ukara kaya mangkono bisa. Ora ana
himpunan !!{sentence}s logika orde kapisan murni kang dicukupi
ing !!a{structure} yen lan mung yen domaine winates.

\begin{prop}
$\Sat{M}{\fn{Inf}}$ yen lan mung yen $\Domain{M}$ tanpa wates.
\end{prop}

\begin{proof}
$\Sat{M}{\fn{Inf}}$ yen lan mung yen
  $\Sat{M}{\lforall[x][\lforall[y][(\eq[u(x)][u(y)] \lif \eq[x][y])]]
    \land \lexists[y][\lforall[x][\eq/[y][u(x)]]]}[s]$
  kanggo sawijining~$s$. Yen mangkono, $s(u)$ iku fungsi
  !!a{injective}, lan ana $y \in \Domain{M}$ kang ora
  kalebu jangkauane~$s(u)$. Kosok baline, yen ana
  !!a{injective} $f\colon \Domain{M} \to \Domain{M}$
  kanthi $\ran{f} \neq \Domain{M}$, mula $s(u) = f$
  minangka pangaji variabel kang dibutuhake.
\end{proof}

Himpunan $M$ kang ora kosong iku !!{enumerable} yen ana enumerasi
\[
m_0, m_1, m_2, \dots
\]
saka kabeh !!{element}s ing himpunan iku (tanpa pangulangan nanging bisa
winates). Kanggo himpunan ora kosong, enumerasi kaya mangkono
ana yen lan mung yen ana !!a{element} $z \in M$ lan fungsi
$f\colon M \to M$ saengga $z$, $f(z)$, $f(f(z))$, \dots,
iku kabeh !!{element}s saka~$M$. Yen enumerasi kasebut ana,
$z = m_0$ lan $f(m_k) = m_{k+1}$ (utawa $f(m_k) = m_k$
yen $m_k$ iku !!{element} pungkasan enumerasi) iku
!!{element} lan fungsi kang dibutuhake. Kosok baline, yen
$z$ lan $f$ kaya mangkono ana, mula $z$, $f(z)$,
$f(f(z))$, \dots, iku enumerasi saka~$M$, lan $M$
!!{enumerable}. Miturut definisi sadurunge, himpunan kosong
bisa dienumerasi nganggo dhaptar kosong; domain sawijining
struktur tansah ora kosong. Cathetan koreksi sumber OLPL-804: definisi enumerasi bisa nyakup himpunan kosong, nanging gambaran nganggo unsur wiwitan z mung lumaku kanggo himpunan ora kosong; domain struktur ing kene pancen ora kosong. Anane $z$ lan~$f$ bisa
dinyatakake ing logika orde kapindho kanggo ngasilake
!!{sentence} kang bener ing !!a{structure} yen lan mung
yen !!{structure} kasebut !!{enumerable}:
\[
\fn{Count} \ident
\lexists[z][\lexists[u][\lforall[X][((X(z) \land
      \lforall[x][(X(x) \lif X(u(x)))]) \lif \lforall[x][X(x)])]]]
\]

\begin{prop}
$\Sat{M}{\fn{Count}}$ yen lan mung yen $\Domain{M}$
!!{enumerable}.
\end{prop}

\begin{proof}
Upamana $\Domain{M}$ iku !!{enumerable}, lan upamana
$m_0$, $m_1$, \dots, sawijining enumerasi. Kanthi
mbusak pangulangan, kita bisa njamin yen ora ana $m_k$
kang katon kaping pindho. Tetepna $f(m_k) = m_{k+1}$
lan $s(z) = m_0$ sarta $s(u) = f$. Yen enumerasine
winates, fungsi kasebut ngirim unsur pungkasan menyang
awake dhewe, kaya kang diterangake ing ndhuwur. Cathetan koreksi sumber OLPL-805: kanggo enumerasi winates, aturan f ing unsur pungkasan perlu ditetepake, umpamane ngirim unsur iku menyang awake dhewe, supaya f tetep fungsi total. Kita bakal
nuduhake yen
\[
\Sat{M}{\lforall[X][((X(z) \land \lforall[x][(X(x) \lif X(u(x)))])
    \lif \lforall[x][X(x)])]}[s]
\]
Upamana $M \subseteq \Domain{M}$ sembarang. Upamana uga
$\Sat{M}{(X(z) \land \lforall[x][(X(x) \lif
X(u(x)))])}[\Subst{s}{M}{X}]$. Mula
$\Subst{s}{M}{X}(z) \in M$, lan yen $x \in M$, uga
$(\Subst{s}{M}{X}(u))(x) \in M$. Kanthi tembung liya,
amarga $\varAssign{\Subst{s}{M}{X}}{s}{X}$,
$m_0 \in M$ lan yen $x \in M$ mula $f(x) \in M$,
dadi $m_0 \in M$, $m_1 = f(m_0) \in M$,
$m_2 = f(f(m_0)) \in M$, lan sateruse.
Mula $M = \Domain{M}$, saengga
$\Sat{M}{\lforall[x][X(x)]}[\Subst{s}{M}{X}]$.
Amarga $M \subseteq \Domain{M}$ dipilih sembarang, bukti
iki rampung:
$\Sat{M}{\fn{Count}}$.

Saiki upamana $\Sat{M}{\fn{Count}}$, yaiku
\[
\Sat{M}{\lforall[X][((X(z) \land \lforall[x][(X(x) \lif X(u(x)))])
    \lif \lforall[x][X(x)])]}[s]
\]
kanggo sawijining~$s$. Upamana $m = s(z)$ lan
$f = s(u)$, banjur gatekna $M = \{m, f(m), f(f(m)), \dots\}$.
$M$ kang ditetepake mangkene cetha !!{enumerable}.
Mula
\[
\Sat{M}{(X(z) \land \lforall[x][(X(x) \lif X(u(x)))])
    \lif \lforall[x][X(x)]}[\Subst{s}{M}{X}]
\]
miturut anggepan. Kajaba iku,
$\Sat{M}{X(z)}[\Subst{s}{M}{X}]$ amarga
$M \ni m = \Subst{s}{M}{X}(z)$, lan uga
$\Sat{M}{\lforall[x][(X(x) \lif
X(u(x)))]}[\Subst{s}{M}{X}]$ amarga yen $x \in M$,
uga $f(x) \in M$. Dadi, amarga anteseden lan implikasine
loro-lorone dicukupi, konsekuene uga kudu dicukupi:
$\Sat{M}{\lforall[x][X(x)]}[\Subst{s}{M}{X}]$.
Iki ateges $M = \Domain{M}$, mula $\Domain{M}$
!!{enumerable} amarga $M$ mangkono miturut definisi.
\end{proof}

\begin{prob}
!!{sentence} $\fn{Inf} \land \fn{Count}$ bener persis
ing kabeh domain kang !!{denumerable}. Owahana definisi
$\fn{Count}$ supaya dadi !!{sentence} liya kang langsung
nyatakake yen domaine !!{denumerable}, lan buktikna.
\end{prob}

\end{document}

